\chapter{How Unsupervised Learning Works}
\label{ch:unsupervised}

\chapterguide%
  {The shared foundations in Stages 1--2 and the supervised workflow in Stage 3A for comparison.}
  {explain what changes when no target is supplied, distinguish the main unsupervised tasks, evaluate structure without pretending there is always a ground-truth answer, and follow a safe unsupervised workflow.}

Stage 4A begins by resetting the workflow before any unsupervised model family is chosen. Until now, the main modeling chapters had a target $y$. That target gave training a direction and made evaluation comparatively direct. Unsupervised learning removes that answer key. The input is usually just a feature matrix $X$, and the goal is to discover or construct useful structure.

\begin{mlfigure}
\begin{tikzpicture}[
  bx/.style={draw=MLBlue,fill=MLPaleBlue,rounded corners=2mm,text width=4.4cm,minimum height=1.2cm,align=center,font=\small\bfseries\color{MLNavy}},
  arr/.style={-{Stealth[length=2.2mm]},thick,draw=MLTeal}
]
\node[bx] (s) at (-4.3,0) {Supervised\\$X+y \longrightarrow$ predict $y$};
\node[bx] (u) at (4.3,0) {Unsupervised\\$X \longrightarrow$ discover structure};
\draw[arr] (s.east) -- node[above,font=\scriptsize\color{MLMuted}]{remove the supplied target} (u.west);
\end{tikzpicture}
\end{mlfigure}

\begin{firstread}
Remember one question: \textbf{what structure am I trying to preserve or discover?} ``Unsupervised learning'' is not one objective. Clustering asks for groups, dimensionality reduction asks for a useful lower-dimensional representation, anomaly detection asks which cases are unusual, and density estimation asks where probability mass lies. Association / pattern mining asks which items repeatedly occur together; it appears on the field map for orientation but is \textbf{beyond the core of this edition}.
\end{firstread}

\section{The unsupervised-learning map}

\begin{mltable}
\begin{tabularx}{\linewidth}{@{}>{\bfseries\leavevmode\color{MLNavy}}p{0.27\linewidth}p{0.27\linewidth}X@{}}
\toprule
\rowcolor{MLPaleBlue!92}
Question & First framing & Main output \\
\midrule
Do observations form groups? & Clustering & Cluster assignments, centers, hierarchy, or soft memberships \\
Can the data be represented with fewer dimensions? & Dimensionality reduction / representation & Components or lower-dimensional coordinates \\
Which observations are unusual? & Anomaly / novelty detection & Anomaly score or ranked suspicious cases \\
Where is the data concentrated? & Density estimation & Estimated probability density or mixture components \\
Which items occur together? & Association / pattern mining \textit{(beyond core)} & Frequent itemsets and rules \\
\bottomrule
\end{tabularx}
\end{mltable}

\begin{keyidea}
These tasks can share algorithms and geometry without sharing an evaluation rule. A good clustering is not defined by explained variance; a good PCA representation is not defined by silhouette score; a useful anomaly detector may be judged by review precision rather than a cluster metric.
\end{keyidea}

\section{What does fitting mean without labels?}

An unsupervised method still has an objective; it simply does not compare predictions with a supplied target during fitting.

\begin{mltable}
\begin{tabularx}{\linewidth}{@{}>{\bfseries\leavevmode\color{MLNavy}}p{0.24\linewidth}p{0.34\linewidth}X@{}}
\toprule
\rowcolor{MLPaleBlue!92}
Method family & What fitting tries to do & What is learned \\
\midrule
$K$-means & Reduce squared distance from points to assigned centers & Cluster centers and assignments \\
PCA & Preserve as much variance as possible in a lower-dimensional linear representation & Principal directions and projected coordinates \\
Density / mixture models & Assign high likelihood to the observed feature distribution under a chosen family & Distribution parameters and soft memberships \\
Isolation / local methods & Make unusual observations easy to isolate or locally sparse & Anomaly scores / rankings \\
Association mining \textit{(beyond core)} & Find item combinations that occur often enough to be interesting & Frequent itemsets and conditional rules \\
\bottomrule
\end{tabularx}
\end{mltable}

\begin{optionaltopic}
Association / pattern mining is included here so the unsupervised map is complete, but it is not developed as a Stage 4B model family in this beginner edition. Stage 4B deliberately concentrates on three reusable objectives: grouping observations, building representations, and ranking unusual cases.
\end{optionaltopic}

\section{How do we evaluate when there is no answer key?}

There is no universal ``unsupervised accuracy.'' Use evidence that matches the task.

\subsection{Internal evidence}
Uses the features and fitted structure itself: cluster compactness/separation, reconstruction error, explained variance, density, or stability under resampling and reasonable perturbations.

\subsection{External evidence}
Sometimes known labels exist for analysis even though the algorithm did not use them for fitting. Metrics such as adjusted Rand index can compare a clustering with those labels. Treat this as an external diagnostic, not proof that the discovered clusters are the only valid structure.

\subsection{Practical or downstream evidence}
Often the strongest question is whether the representation or discovered structure is useful: are clusters interpretable and stable, does a compressed representation preserve a downstream signal, and do highly ranked anomalies correspond to cases a reviewer cares about?

\begin{pitfall}
Do not evaluate an unsupervised method by inventing a convenient target after the fact and then treating that target as ground truth. Clusters may cut across known categories for legitimate reasons, and a visually pleasing two-dimensional plot can distort neighborhoods. State what property the method was designed to preserve.
\end{pitfall}

\section{The unsupervised workflow}

\begin{mlfigure}
\begin{tikzpicture}[
  bx/.style={draw=MLBlue,fill=MLPaleBlue,rounded corners=1.5mm,text width=2.0cm,inner xsep=2pt,minimum height=0.9cm,align=center,font=\scriptsize\bfseries\color{MLNavy}},
  arr/.style={-{Stealth[length=2mm]},thick,draw=MLTeal}
]
\node[bx] (a) at (-6.5,0) {1. Define the\\structure sought};
\node[bx] (b) at (-3.9,0) {2. Choose the\\unit and features};
\node[bx] (c) at (-1.3,0) {3. Prepare and\\scale appropriately};
\node[bx] (d) at (1.3,0) {4. Fit candidate\\structures};
\node[bx] (e) at (3.9,0) {5. Check stability\\and meaning};
\node[bx] (f) at (6.5,0) {6. Use on new or\\downstream data};
\draw[arr] (a)--(b); \draw[arr] (b)--(c); \draw[arr] (c)--(d); \draw[arr] (d)--(e); \draw[arr] (e)--(f);
\end{tikzpicture}
\end{mlfigure}

A train/test split is not automatically required for every exploratory clustering of one fixed dataset. It \emph{is} required when you make a generalization claim, tune a reusable transform using downstream labels, build a novelty detector for future observations, or otherwise claim that the learned structure will transfer. The information-boundary rule from \chapref{ch:data} still applies whenever a transformation is fitted and then evaluated on held-out data.

\section{Scaling and representation matter even more here}

Many unsupervised algorithms see only geometry. If annual income is measured in hundreds of thousands and age in tens, a raw Euclidean distance can become mostly an income distance. Standardization, feature weighting, encoding, and the choice of distance therefore define part of the problem rather than merely polishing the input.

\begin{keyidea}
In supervised learning, a target can sometimes reveal that a poor representation is unhelpful. In unsupervised learning there may be no such corrective signal. The representation and distance choice therefore deserve explicit justification before interpreting the structure that appears.
\end{keyidea}

\section{The sequence ahead}

\begin{description}[leftmargin=1.6in,style=nextline]
  \item[\chapref{ch:clustering}] asks whether observations form meaningful groups and compares center-, hierarchy-, density-, and probability-based approaches.
  \item[\chapref{ch:dimred}] asks what information should be preserved when features are compressed, with PCA as the core method and nonlinear visualization as an extension.
  \item[\chapref{ch:anomaly}] asks which observations are unusual and how anomaly scores become operational alerts.
\end{description}

\section{Kaggle implementation: an unsupervised workflow without labels}
\label{sec:kaggle-mall-unsupervised}

\begin{kaggleproject}{Mall Customers --- scaling, repeated clustering, and internal evidence}
\kagglesource{https://www.kaggle.com/datasets/vjchoudhary7/customer-segmentation-tutorial-in-python}{Mall Customer Segmentation Data}
\textbf{Question.} Without a target column, how can we compare candidate values of $k$ using compactness and silhouette structure?
\end{kaggleproject}

The script scales age, annual income, and spending score, then reports K-means inertia and silhouette score for several candidate cluster counts. Neither quantity is a ground-truth accuracy measure; both are internal evidence about the representation and partition.

\textbf{Before running.} Predict the qualitative pattern you expect from the experiment before you look at the output or plot.

\begin{labmode}
Stage 4A is about evidence without labels. Decide what would count as useful evidence before looking at the reported inertia and silhouette values.
\end{labmode}

\lstinputlisting[style=mlpython]{all_experiments/lab14-mall_unsupervised_workflow.py}


\begin{labdeliverable}
Submit the scaled feature definition, inertia and silhouette evidence across candidate $k$, your chosen candidate partition, one visualization, and a decision memo stating what additional domain or stability evidence is required before treating the groups as meaningful customer segments.
\end{labdeliverable}



\begin{quickcheck}
\begin{enumerate}
  \item Why is there no universal unsupervised-learning accuracy score?
  \item Give one internal, one external, and one downstream form of evidence that could support an unsupervised result.
  \item When does an unsupervised project need a held-out split even though no target is used during fitting?
\end{enumerate}
\end{quickcheck}

\begin{chapterrecap}
\begin{itemize}
  \item Unsupervised learning fits structure in $X$ without using a supplied target during fitting.
  \item Clustering, dimensionality reduction, anomaly detection, and density estimation answer different questions and therefore need different objectives and evaluation evidence.
  \item Internal, external, and practical/downstream evaluation are complementary; there is no universal unsupervised accuracy score.
  \item Scaling, representation, and distance choices are part of the model because many unsupervised methods are geometric.
\end{itemize}
\end{chapterrecap}
